\begin{explanationbox}
	\textbf{\textcolor{CExplanation}{一句话直觉}}

	当一组平行线同时穿过两条直线时，它们被截得的线段长度比例保持一致。

	\medskip
	\textbf{\textcolor{CExplanation}{核心拆解}}
	\begin{description}[style=nextline, leftmargin=2.8em, labelsep=0.5em, itemsep=0.45em]
		\item[\textbf{要点一（三线平行组）：}] 至少三条不重合的平行线，如 $l_1\parallel l_2\parallel l_3$，它们之间的距离可以不相等。
		\item[\textbf{要点二（对应次序）：}] 截线的交点按平行线顺序对应，$A\leftrightarrow D$、$B\leftrightarrow E$、$C\leftrightarrow F$，从而 $AB$ 对应 $DE$，$BC$ 对应 $EF$，$AC$ 对应 $DF$。
		\item[\textbf{要点三（比例恒等）：}] 无论两条截线是否相交，对应线段的比值总是相等的，即 $\dfrac{AB}{BC}=\dfrac{DE}{EF}$ 和 $\dfrac{AB}{AC}=\dfrac{DE}{DF}$。
	\end{description}

	\medskip
	\textbf{\textcolor{CExplanation}{几何本质（平行构相似）}}

	通过过点作平行线，可以构造相似三角形；或者在截线相交时直接由同位角相等得到三角形相似，从而将线段比转化为相似比。

	\medskip
	\textbf{\textcolor{CExplanation}{代数意义（分数等式）}}

	比例式 $\dfrac{AB}{BC}=\dfrac{DE}{EF}$ 可转化为 $AB\cdot EF = BC\cdot DE$，便于解出未知线段长度。

	\medskip
	\textbf{\textcolor{CExplanation}{考点价值（比例工具）}}
	\begin{itemize}[leftmargin=*, itemsep=0.5em, topsep=0.4em]
		\item \textbf{考法一（求线段长）：} 已知部分线段长和比例关系，求未知线段。
		\item \textbf{考法二（证比例式）：} 在几何证明中，通过构造平行线组推导比例关系。
	\end{itemize}

	\medskip
	\textbf{\textcolor{CExplanation}{顿悟点}}

	找出哪两条是截线，哪一组是平行线，然后严格按照交点顺序写出比例式。

	\medskip
	\textbf{\textcolor{CExplanation}{使用场景}}
	\begin{itemize}[leftmargin=*, itemsep=0.5em, topsep=0.4em]
		\item \textbf{场景一（梯形问题）：} 梯形中作与两底平行的直线，求线段长度或证明比例。
		\item \textbf{场景二（三角形截线）：} 在三角形中作平行于一边的直线，直接使用截线段成比例。
	\end{itemize}
\end{explanationbox}
